\begin{textsample}{POS Dim 3 – pr – Score 8.00 – pr/pr\_inf\_ef\_10.txt}  \label{ex:f3_pos_259}
Transição entre as etapas da educação básica Em 2005, pela Lei Federal n.º 11.114/05 ( BRASIL, 2005 ), que alterou o Artigo 6.º da LDBEN, tornou -se obrigatória a matrícula da \textbf{criança} aos seis anos de idade no ensino fundamental, sendo o mesmo ampliado para nove anos de duração por meio da Lei n.º 11.274/2006 ( BRASIL, 2006 ).
Essa legislação atendeu ao disposto no Plano Nacional de Educação de 2001, Lei n.º 10.172/2001 ( BRASIL, 2001 ), que, entre suas metas, estabeleceu tal ampliação.
Posteriormente, a ampliação da obrigatoriedade da educação no Brasil passou a ser dos 4 aos 17 anos de idade pela Emenda Constitucional n.º 59/2009 ( BRASIL, 2009 ), regulamentada pela Lei n.º 12.796/2013, estendendo a obrigatoriedade da escolarização às etapas da Educação \textbf{Infantil} ( Pré-escola ) e ao Ensino Médio, alterando o artigo 4º da LDBEN.
Diante dos direitos de aprendizagens dispostos no texto da Base Nacional Comum Curricular, fica explícito que todos os estudantes devem ter as mesmas oportunidades de aprendizagem.
Isto posto, a escolarização da infância, ou seja, da Educação \textbf{Infantil} ao Ensino Fundamental - Anos Iniciais, deve ser estabelecida por práticas educativas específicas visando ao desenvolvimento e à aprendizagem das \textbf{crianças} em suas diferentes faixas etárias e processos formativos.
Portanto, os tempos e espaços devem ser diferenciados, posicionando os estudantes em lugares distintos.
A transição entre a Educação \textbf{Infantil} e o Ensino Fundamental é um momento crucial e complexo na vida das \textbf{crianças} e as instituições de ensino devem constituir ações que minimizem a ruptura que pode ser causada.
O primordial é ter como critério que a educação \textbf{infantil} não se ocupa da preparação para a entrada no ensino fundamental, mas que, em cada ação e prática, o movimento seja de atender às especificidades, individualidades e as totalidades das \textbf{crianças}.
Os docentes, sujeitos diretos de contato com os estudantes, devem considerar que a perspectiva formativa nessa etapa se dá por meio do jogo, do \textbf{brinquedo} e da ludicidade.
Neste contexto, é necessário ponderar atentamente para algumas questões que podem nortear as ações finais da educação \textbf{infantil} e iniciais do ensino fundamental : o que significa atender as especificidades da infância?
Quais fatores interferem no processo de transição da educação \textbf{infantil} para o ensino fundamental?
Como trabalhar o"abandono"simbólico dos \textbf{colegas} e referenciais anteriores?
O que implica considerar aspectos que vão para além da adaptação física e estrutural?
Como priorizar a iniciação em conceitos mais complexos?
Como ajudar as \textbf{crianças} a reelaborar afinidades com os professores?
Como organizar e distribuir o espaço de sala de aula e os demais espaços da instituição de ensino em prol das \textbf{crianças}?
Qual o melhor acolhimento às \textbf{crianças} de seis anos no ensino fundamental?
Dessas reflexões surge a necessidade de repensar as práticas pedagógicas relacionadas ao Ensino Fundamental para as \textbf{crianças} que, atualmente, ingressam mais cedo nas escolas : o que prever para a alegria de permanecer nesse espaço?
Como possibilitar a integração e pertencimento da \textbf{criança} nesse novo espaço escolar?
Como favorecer as interações e trocas que possibilitam a aprendizagem das \textbf{crianças}?
O que deve ser avaliado sobre as \textbf{crianças}?
A proposta pedagógica está a favor da \textbf{criança} ou do \textbf{adulto}?
O \textbf{adulto} consegue perceber como a \textbf{criança} aprende?
O que é necessário para melhorar as condições de equidade de aprendizagens e qualidade do ensino?
O que prever de aprendizagens para a alfabetização e o letramento?
Torna -se essencial compreender que a \textbf{criança} advinda da Educação \textbf{Infantil}, com cinco ou seis anos, ainda será \textbf{criança} até os nove ou dez anos de idade.
Respeitar essa etapa da vida humana deve ser o objetivo de trabalho dos docentes e gestores de educação com vistas à formação integral.
Assim, considerando que a educação \textbf{infantil} tem como finalidade atender as \textbf{crianças} em suas especificidades, o uso das linguagens da infância como a \textbf{brincadeira}, o jogo, o faz de conta, a liberdade de pensamento, deve ser mediada pelo docente do ensino fundamental ampliando ou reelaborando as práticas pedagógicas de forma a serem mais coerentes para e com as \textbf{crianças}.
Cada momento de ingresso numa instituição de ensino deve ser organizado com vistas às necessidades físicas, cognitivas e emocionais das \textbf{crianças}, respeitando seus medos e inseguranças, amenizando angústias de adaptação.
O processo de municipalização da oferta do ensino fundamental no Brasil foi intenso ao longo das últimas décadas.
Esse fato ocorreu de forma gradativa e diversa entre os 26 estados da federação, separando em diferentes esferas administrativas, em maior ou \textbf{menor} grau, a fase dos anos iniciais ( 1º ao 5º ano ), que ficou sob a responsabilidade dos municípios, e a fase dos anos finais ( 6º ao 9º ano ), que ficou sob a responsabilidade dos estados.
No Paraná, atualmente, o resultado desse processo significa que a municipalização da oferta dos anos iniciais do ensino fundamental nas escolas públicas chega a 99,49 \% ( BRASIL, 2017 ).
Uma exploração da história sobre como se configurou o ensino fundamental como etapa de educação básica tal como estabelecida atualmente pela LDBEN 9.394/1996 mostra que, a partir da Lei nº 5.692/71 ( BRASIL, 1971 ), que fixou diretrizes e bases para o ensino de 1º e 2º graus, ficou estabelecido o ensino de 1º grau obrigatório dos 7 aos 14 anos.
Diferentemente do que estava prescrito na LDBEN nº 4.024/1961 ( BRASIL, 1961 ), em que essa obrigatoriedade se limitava às quatro séries iniciais do então chamado ensino primário e incluía a dependência de aprovação em exame de admissão para o ingresso no ciclo dos quatro anos seguintes, chamado de ginasial.
Desde a instituição de uma etapa do ensino que agrupou duas organizações pedagógicas diferentes no ensino fundamental obrigatório, sem definir a necessária metodologia articuladora das questões pedagógicas características dessa transição, permaneceu a fragilidade na adequação metodológica, na integração curricular, na correspondente formação de professores, no reconhecimento das diferentes culturas escolares, na integração entre as redes de ensino de modo a articular as informações sobre o desenvolvimento e aprendizagem dos estudantes, na atenção à transição da infância para a adolescência, entre outras articulações.
Nesse cenário, a essencial tarefa organizadora e unificadora do currículo por meio da Base Nacional Comum Curricular, como potencial articulador do ensino fundamental, não se realiza por si só.
É necessário ponderar o indispensável trabalho conjunto de professores, sujeitos que atribuem vitalidade ao currículo e que atuam nas duas fases dessa etapa, de forma que os esforços por conhecer a organização curricular nos anos iniciais e finais, bem como o estabelecimento de estratégias de atuação nessa transição tenham início nos primeiros anos e continuem ocorrendo do 6º ano em diante.
Ou seja, se faz necessária uma atenção especial na reflexão e viabilização de práticas pedagógicas que integrem os envolvidos no processo, tendo como elemento indutor uma política educacional articuladora entre as etapas e fases : da \textbf{creche} para \textbf{pré-escola}, da \textbf{pré-escola} para os anos iniciais do ensino fundamental e destes para os anos finais.
Esse esforço de ampliação das oportunidades de sucesso do estudante pode possibilitar efetivamente o desenvolvimento integral do estudante.

% matched lemmas: adulto, brincadeira, brinquedo, colega, creche, criança, infantil, pequeno, pré-escola
\end{textsample}
